\section{引言}

全切片图像（whole-slide images，WSIs）以极高分辨率记录完整的组织形态，是计算病理学和计算机辅助临床分析的重要基础。准确的 WSI 组织分割可支持肿瘤面积与组织比例测量、肿瘤微环境分析、空间生物标志物研究、感兴趣区域筛选，以及下游诊断或预后建模。然而，单张 WSI 可能包含数十亿像素，同时伴有复杂的组织边界和显著的尺度变化。因此，由病理学家逐像素勾画组织区域不仅成本高昂，也难以扩展至大规模数据。在有限人工投入下获得准确且可灵活适配的组织图谱，仍是一项亟待解决的重要问题。

现有方法大多将组织分割建模为预定义类别集合上的多分类预测，例如肿瘤、间质、坏死和脂肪组织。当类别定义稳定且具备密集标注时，这一范式能够有效工作，但其输出空间仅限于训练阶段见过的类别。对于新的组织结构、具有特定细胞组成的区域，或面向某项具体研究定义的目标，通常需要重新获取像素级标注并训练模型。此外，不同癌种、患者队列和临床问题对组织的定义并不一致，这使得固定类别模型难以作为可扩展的组织分析工具。

以分割一切模型（Segment Anything Model，SAM）为代表的可提示分割模型提供了另一种路径：用户通过少量点、边界框或涂画来指定目标，模型无需重新训练特定类别的预测头即可输出相应掩膜。然而，将此类模型直接应用于 WSI 仍面临三方面困难。首先，组织身份并不总能仅凭颜色、纹理或局部边界加以区分，它还可能取决于细胞类型、细胞密度及其空间组成。其次，用户可能只在局部标注某类组织的一个实例，而该组织往往分散于切片中多个彼此不连通的位置。主要依赖邻近区域外观的模型，可能无法将用户意图所对应的组织语义传播至整张切片。最后，病理组织同时涉及细胞形态、局部组织结构与全局空间布局，单一尺度通常难以兼顾精细边界定位与可靠的组织身份判别。

因此，提示驱动的 WSI 分割不应局限于围绕点击位置扩展边界。更具挑战性的核心问题在于：模型能否从一次或少数几次交互中推断用户意图所指的组织概念，并在十亿像素级切片中检索所有语义一致的区域。解决这一问题，需要构建同时表征视觉外观与细胞组成的特征表示，并具备跨组织层级进行推理的机制。

我们的核心观察是，组织由具有特定身份、密度和空间分布模式的细胞构成。低倍视野下外观相似的区域，其细胞组成可能并不相同；相反，同类组织中彼此不连通的实例往往共享相似的细胞与组织结构模式。基于这一观察，我们将组织表示为可学习区域，使其同时融合视觉轮廓与细胞证据，并在精细、中间和粗粒度三个尺度间交换信息。

图~\ref{fig:overview} 概述了 FewClick：一个面向 WSI 少提示组织分割的细胞感知层次化区域提示框架。该框架提取空间对齐的多尺度视觉特征，并通过可微分软分配形成可学习的组织区域；随后利用细胞到区域注意力，将细胞嵌入、密度及组成信息注入区域表示。稀疏的父子层次结构与门控适配器进一步融合局部边界、邻域信息和宏观上下文。在交互阶段，两个相互独立且具有置换不变性的编码器分别聚合正提示与负提示。上下文感知解码器将所得目标表示与整张 WSI 中的所有区域进行匹配，再将区域概率投影回像素空间。与局部连通区域扩展不同，该框架能够检索在空间上彼此不连通、但具有一致细胞组成和组织语义的区域。

本文的主要贡献包括四个方面：（1）我们将 WSI 分割表述为由提示定义的二值组织检索问题，而非基于固定输出通道的多分类预测；（2）我们提出可微分的细胞感知组织区域，以统一编码视觉特征、细胞嵌入、细胞密度与细胞组成；（3）我们在空间对齐的精细、中间和粗粒度区域之间设计了稀疏的自底向上与自顶向下交互机制；（4）我们将带符号的提示编码与上下文感知区域解码相结合，使正、负证据均能参与整张切片范围内的推理。受控对比实验表明，FewClick 相较三种可提示分割基线均取得了稳定提升；消融实验进一步验证了可学习区域、细胞证据和多尺度上下文各自的贡献。
